\section{Numerical study}\label{sec:numerics}

\begin{rv}{revM}{\tE{7}\tA{N5}}
The computations in this section illustrate the theory for a
single benchmark model. They do not validate the theorems, and they
establish no property of the continuous problem beyond what the refinement studies indicate. 
\end{rv}

\begin{rv}{revM}{\tA{N4}\tB{21, 23}}
\subsection{Benchmark and discretization}

The benchmark is the reflected correlated Ornstein--Uhlenbeck system
\[
\dd X^i_t=\kappa_i(\theta_i-X^i_t)\,\dd t+\sigma_i\,\dd W^i_t+\dd L^i_t,\quad i=1,2,\qquad\dd\langle W^1,W^2\rangle_t=\rho\,\dd t,
\]
with $\kappa_i>0$, so that each coordinate reverts to $\theta_i$, and running cost
$c(x)=c_0+c_1x_1+c_2x_2$. The covariance is constant and positive definite for $|\rho|<1$, so
Assumption~\ref{ass:state} holds (Remark~\ref{rem:vertex}), and Assumption~\ref{ass:growth} holds with $m=1$
(Remark~\ref{rem:lyap}). Here $q=\sigma_1^2-2\alpha\rho\sigma_1\sigma_2+\alpha^2\sigma_2^2$, with
$\partial q/\partial\rho=-2\alpha\sigma_1\sigma_2<0$. Baseline parameters and controls are listed in
\Cref{tab:params}.

\begin{table}[htbp]
\begin{rv}{revA}{\tA{N4}}
\centering\small
\caption{Baseline parameters and computational controls.}\label{tab:params}
\begin{tabular}{@{}lll@{}}
\toprule
Symbol & Value & Meaning\\
\midrule
$r$ & $0.05$ & discount rate\\
$\alpha$ & $1.0$ & relative value of alternative 2\\
$\kappa_1,\kappa_2$ & $0.4,\ 0.4$ & mean-reversion speeds\\
$\theta_1,\theta_2$ & $1.0,\ 1.0$ & long-run levels\\
$\sigma_1,\sigma_2$ & $0.4,\ 0.4$ & volatilities\\
$\rho$ & $0.0$ & correlation\\
$c_0,c_1,c_2$ & $0.01,\ 0,\ 0$ & running-cost parameters\\
\midrule
$X_{\max}$; $N$ & $3$ (and $4$); $61$--$481$ & FD domain; grid points per axis\\
window & $[0,2.5]^2$ & areas and loss statistics (uniform node weights)\\
$x_{\rm ref}$ & $(0.9,1.0)$ & reference state for the Monte Carlo study\\
$\Delta t$; $T_{\max}$; $M$ & $0.0025$--$0.02$; $40$; $2\times10^5$ & MC step, horizon, paths (\Cref{sec:mcnum})\\
random numbers & PCG64, seed $20261005$ & independent paths (\Cref{sec:mcnum,sec:biasnum})\\
\bottomrule
\end{tabular}
\end{rv}
\end{table}

At the baseline all drift differences are central, so the scheme of \Cref{sec:fd} is second-order
consistent in the interior and monotone; for the correlated runs $a_{ii}-|a_{12}|=\sigma^2(1-|\rho|)>0$
and the seven-point stencil is monotone. Howard's iteration terminated in $15$, $28$, $54$ and $107$
steps for $N=61,121,241,481$. Free boundaries are located inside grid cells as follows: near a
smooth-fit boundary $U-G$ behaves like a constant times the squared distance to the boundary, so
$\sqrt{U-G}$ is extrapolated linearly from the last two continuation nodes to its zero. Areas and
loss statistics are computed over the nodes of the window $[0,2.5]^2$, each node carrying weight
$h^2$.
\end{rv}

\begin{rv}{revM}{\tE{2, 7}\tA{N4, Ed2}\tB{22}}
\subsection{Stopping set, refinement, and the distance to the indifference line}\label{sec:geomnum}

\Cref{fig:regions} shows the computed stopping set for $N=481$. For this benchmark it consists of
two arms, one in $E_1$ and one in $E_2$, separated by a continuation region that contains $\Da$
and narrows as $x_1=x_2$ increases. This is an observation about the computed solution of one
model. It is consistent with \Cref{thm:diag} and Corollary~\ref{cor:geometry}, which it does not prove and
which do not predict it: by Remark~\ref{rem:geometries}, the theorems would be equally consistent with a
one-sided or empty stopping set. Corollary~\ref{cor:geometry}(b) gives the checkable necessary condition
$\cD\cap E_1\cap\Q\subset\{\Gac\ge0\}=\{x_1\ge(\kappa_1\theta_1-c_0)/(r+\kappa_1)\}=\{x_1\ge0.867\}$;
the computed arm lies in $\{x_1\ge1.46\}$.

\begin{figure}[htbp]
\begin{rv}{revM}{\tE{2}\tA{N4}}
\centering
\includegraphics[width=0.52\linewidth]{figs/regions.pdf}
\caption{Computed stopping set for the baseline ($N=481$): arms $\cD^{(1)}$ and $\cD^{(2)}$, the
indifference line $\Da$ (dashed), the necessary region boundary $\Gac=0$ of
Corollary~\ref{cor:geometry}(b) (dash-dotted), and the formal wedge half-width $q/(4\Gac)$ of
Remark~\ref{rem:width} (dotted).}\label{fig:regions}
\end{rv}
\end{figure}

\Cref{tab:refine} and \Cref{fig:refine} report the refinement study. The value at $x_{\rm ref}$
is stable to four digits and the boundary positions inside the window to within about $0.01$; both
are insensitive to the truncation $X_{\max}\in\{3,4\}$. \Cref{tab:refine} shows that the minimum nodal distance over the full domain equals
$h$ at every resolution. It is attained in the outer part of the domain and is produced by the reflecting wall at
$X_{\max}$, at which $G$ does not satisfy the Neumann condition. The computations resolve the wedge inside the window: its half-width decreases from $0.55$ at $x_2=1$ to $0.105$ at $x_2=2.5$,
within $3$--$11\%$ of the formal value $q/(4\Gac)$ for $x_2\ge1.5$, and the continuation premium on
$\Da$ at $(2.5,2.5)$ is $0.0274$ against the formal $q/(16\Gac)=0.0272$. Since $\Gac\to\infty$
along $\Da$, both tend to zero: the computations support the picture of Remark~\ref{rem:geometries}(iii),
strict continuation on $\Da$ with no uniform separation.
\end{rv}

\begin{table}[htbp]
\begin{rv}{revM}{\tA{N4}\tB{22}}
\centering\small
\caption{Grid refinement for the baseline, $X_{\max}=3$. The runs with $X_{\max}=4$ and the same
$h\ge0.0125$ agree with the corresponding rows. The last line gives the formal
prediction $q/(4\Gac)$ of Remark~\ref{rem:width}.}\label{tab:refine}
\begin{tabular}{@{}rrcccccc@{}}
\toprule
 & & & & \multicolumn{3}{c}{half-width $b_1(x_2)-x_2$ at $x_2=$} & nodal min.\ distance\\
\cmidrule(lr){5-7}
$N$ & $h$ & $(V-G)(x_{\rm ref})$ & $b_1(0)$ & $1.5$ & $2.0$ & $2.5$ & to $\Da$, full domain\\
\midrule
$61$ & $0.05$ & $0.27529$ & $1.466$ & $0.249$ & $0.148$ & $0.106$ & $0.0500$\\
$121$ & $0.025$ & $0.27521$ & $1.457$ & $0.249$ & $0.147$ & $0.105$ & $0.0250$\\
$241$ & $0.0125$ & $0.27519$ & $1.463$ & $0.249$ & $0.147$ & $0.104$ & $0.0125$\\
$481$ & $0.00625$ & $0.27518$ & $1.463$ & $0.250$ & $0.146$ & $0.105$ & $0.0063$\\
\midrule
\multicolumn{4}{@{}l}{formal $q/(4\Gac)$} & $0.281$ & $0.157$ & $0.109$ & ---\\
\bottomrule
\end{tabular}
\end{rv}
\end{table}

\begin{figure}[htbp]
\begin{rv}{revM}{\tA{N4}\tB{22}}
\centering
\includegraphics[width=0.95\linewidth]{figs/refine.pdf}
\caption{Left: half-width of the continuation wedge along $\Da$ for three grids, with the formal
prediction. Right: the minimum over the whole computational domain of the nodal distance between
$\cD^{(1)}$ and $\Da$ equals one grid cell at every resolution, whereas the reconstructed distance
within the window $[0,2.5]^2$ is resolved.}\label{fig:refine}
\end{rv}
\end{figure}

\begin{rv}{revM}{\tE{7}\tA{N1}}
\subsection{Consistency diagnostics}

\Cref{tab:residuals} lists the diagnostics of \Cref{sec:diagnostics}. The discrete
variational-inequality residuals are at round-off level, which shows that the discrete problem
\eqref{eq:discreteVI} is solved. The smooth-fit quotient halves with $h$, as it
should when the continuous solution has no derivative jump at the free boundary, and the Neumann
quotient decreases at second order. These are consistency diagnostics. 
\end{rv}

\begin{table}[htbp]
\begin{rv}{revA}{\tA{N1}}
\centering\small
\caption{Consistency diagnostics for the finite-difference solution (\Cref{sec:diagnostics}). The
first four rows vanish up to linear-solver round-off by construction of Howard's
algorithm.}\label{tab:residuals}
\begin{tabular}{@{}lccc@{}}
\toprule
Quantity & $N=121$ & $N=241$ & $N=481$\\
\midrule
majorization $\max(G-U)^+$ & $0$ & $0$ & $0$\\
superharmonicity $\max(-((rI-\cL_h)U+c))^+$ & $5\times10^{-13}$ & $3\times10^{-12}$ & $1\times10^{-11}$\\
complementarity $\max|((rI-\cL_h)U+c)(U-G)|$ & $3\times10^{-13}$ & $1\times10^{-12}$ & $5\times10^{-12}$\\
continuation equation $\max_{\cC_h}|(rI-\cL_h)U+c|$ & $6\times10^{-13}$ & $3\times10^{-12}$ & $1\times10^{-11}$\\
smooth-fit quotient $\max(U-G)(z-he_1)/h$ & $0.110$ & $0.059$ & $0.028$\\
Neumann quotient at the axes & $5.8\times10^{-4}$ & $1.5\times10^{-4}$ & $3.9\times10^{-5}$\\
\bottomrule
\end{tabular}
\end{rv}
\end{table}

\begin{rv}{revM}{\tE{2}\tA{M2, N3}}
\subsection{Comparison policies}\label{sec:policies}

All policies are evaluated under the true dynamics and the true payoff by the linear solve of
\Cref{sec:fd} with $N=241$, and compared with the finite-difference solution through the relative
loss \eqref{eq:loss} over the window. The mean of $\hat V_{\rm ref}$ over the window is $1.769$.

\emph{One-sided policies.} Policy (a) retains the $\cD^{(2)}$ arm of the computed solution and never
stops in $E_1$. It measures the performance of one
particular one-sided policy, not the best that a one-sided class can do. Policy (b) is the optimum
among all policies that are allowed to stop only in $E_2$, obtained by solving
\eqref{eq:discreteVI} with the obstacle imposed on the nodes of $E_2$ only. Its stopping set is an
epigraph $\{x_2\ge b(x_1)\}$; for $x_1\gtrsim1.4$ it extends down to the first row of nodes above
$\Da$, because a policy that cannot exercise alternative~1 stops as soon as it is allowed to. The
loss of (b) is smaller than that of (a) but of the same order (mean $11.7\%$ against $13.6\%$,
\Cref{tab:loss}), and in both cases it is concentrated in the region of $E_1$ where the computed
solution stops (\Cref{fig:loss}, panels a--b). The conclusion is quantitative and specific to this
symmetric benchmark: here, never exercising alternative~1 is costly. We note that one-sided stopping sets are consistent with \Cref{thm:diag} (Remark~\ref{rem:geometries}(i)).
\end{rv}

\begin{table}[htbp]
\begin{rv}{revM}{\tE{2}\tA{N2, N3}}
\centering\small
\caption{Loss of comparison policies relative to the finite-difference solution, over the nodes of
$[0,2.5]^2$ ($N=241$). Relative loss \eqref{eq:loss} in percent; absolute loss
$\hat V_{\rm ref}-\hat J_b$ in value units. (a): only the $\cD^{(2)}$ arm of the computed solution; (b): optimal among policies that stop only in $E_2$; (c): optimal for \eqref{eq:gainform} without the local-time term.}\label{tab:loss}
\begin{tabular}{@{}lrrrrrr@{}}
\toprule
 & \multicolumn{3}{c}{relative loss (\%)} & \multicolumn{3}{c}{absolute loss}\\
\cmidrule(lr){2-4}\cmidrule(lr){5-7}
Policy & mean & 95th pct. & max & mean & 95th pct. & max\\
\midrule
(a) one-sided, retained arm & $13.61$ & $40.91$ & $47.10$ & $0.390$ & $1.345$ & $1.649$\\
(b) one-sided, re-optimised & $11.68$ & $38.65$ & $45.30$ & $0.338$ & $1.268$ & $1.585$\\
(c) local-time reward dropped & $1.91$ & $7.58$ & $13.14$ & $0.043$ & $0.168$ & $0.301$\\
\bottomrule
\end{tabular}
\end{rv}
\end{table}

\begin{figure}[htbp]
\begin{rv}{revM}{\tE{2}\tA{N2, N3}}
\centering
\includegraphics[width=\linewidth]{figs/loss.pdf}
\caption{Relative loss \eqref{eq:loss} of the comparison policies of \Cref{tab:loss} (dashed:
$\Da$; grey: boundary of the computed stopping set; orange: boundary of the comparison policy's
stopping set, where it differs).}\label{fig:loss}
\end{rv}
\end{figure}

\begin{rv}{revA}{\tA{N2}}
\emph{Dropping the local-time reward.} Smoothing changes the reward and spreads the curvature over a
band instead of deleting it. The gain form \eqref{eq:gainform} provides a well-defined ablation
that leaves the reward unchanged. Policy (c) stops on the set $\{\tilde w=0\}$, where
\[
\tilde w(x)=\sup_\tau\E_x\Big[-\int_0^\tau e^{-rs}\Gac(X_s)\,\dd s\Big]
\]
is the value of \eqref{eq:gainform} with the term $\tfrac12\dd\loc$ removed; numerically, $\tilde w$
solves $\min\{(rI-\cL_h)\tilde w+\Gac,\tilde w\}=0$. (With the full discrete gain
$(rI-\cL_h)G+c$ in place of $\Gac$ the same solver reproduces $\hat V_{\rm ref}-G$ to $10^{-14}$,
which checks the discrete analogue of Proposition~\ref{prop:gainform}.) Policy (c) stops on $61\%$ of the
nodes of $\Da$ in the window, namely for $x_1=x_2\gtrsim0.99$, which \Cref{thm:diag} shows to be
suboptimal, and its loss is concentrated around $\Da$ (\Cref{fig:loss}, panel c): mean $1.9\%$,
maximum $13.1\%$. This is the quantitative effect of the singular gain in this benchmark.

\emph{Payoff smoothing.} For completeness, \Cref{tab:smoothing} reports the smoothing experiment
with its definition and a parameter study. The smoothed payoff is
$G_\epsilon=\alpha x_2+\psi_\epsilon(Y)$ with $\psi_\epsilon(y)=\tfrac12(y+\sqrt{y^2+\epsilon^2})$;
its generator has the regular second-order term $\tfrac12q\,\psi_\epsilon''(Y)$ with
$\psi_\epsilon''(y)=\epsilon^2/(2(y^2+\epsilon^2)^{3/2})$, which integrates to one for every
$\epsilon$ and converges to $\delta_0$. If $\|G_\epsilon-G\|_\infty\le\eta$ then
$\|V_\epsilon-V\|_\infty\le\eta$, and a policy optimal for $G_\epsilon$ loses at most $2\eta$ under
$G$; here $\eta=\epsilon/2$. The loss decreases rapidly with $\epsilon$ and is negligible once $\epsilon\le0.1$, where the smoothed-payoff policy no longer stops on $\Da$. Smoothing therefore does not remove the mechanism; a small smoothing parameter reproduces it. 
\end{rv}

\begin{table}[htbp]
\begin{rv}{revA}{\tA{N2}}
\centering\small
\caption{Payoff-smoothing misspecification: policy optimal for
$G_\epsilon=\alpha x_2+\tfrac12\big(Y+\sqrt{Y^2+\epsilon^2}\big)$, evaluated under the true payoff
($N=241$, window $[0,2.5]^2$). Since $0\le G_\epsilon-G\le\epsilon/2$, the absolute loss is at most
$\epsilon$.}\label{tab:smoothing}
\begin{tabular}{@{}rrrrrr@{}}
\toprule
 & \multicolumn{2}{c}{relative loss (\%)} & \multicolumn{2}{c}{absolute loss} & fraction of $\Da$-nodes\\
\cmidrule(lr){2-3}\cmidrule(lr){4-5}
$\epsilon$ & mean & max & mean & max & in the stopping set\\
\midrule
$0.8$ & $0.375$ & $6.32$ & $0.0088$ & $0.153$ & $0.49$\\
$0.4$ & $0.109$ & $3.81$ & $0.0027$ & $0.097$ & $0.42$\\
$0.2$ & $0.027$ & $1.85$ & $0.0008$ & $0.052$ & $0.30$\\
$0.1$ & $0.003$ & $0.31$ & $0.0001$ & $0.011$ & $0$\\
$0.05$ & $2\times10^{-4}$ & $0.017$ & $6\times10^{-6}$ & $6\times10^{-4}$ & $0$\\
$0.025$ & $1\times10^{-5}$ & $0.003$ & $4\times10^{-7}$ & $9\times10^{-5}$ & $0$\\
\bottomrule
\end{tabular}
\end{rv}
\end{table}

\subsection{Signed post-stopping term}\label{sec:biasnum}

\begin{rv}{revA}{\tA{M4, N3, N4}}
By Proposition~\ref{prop:vanish} the signed term $B$
vanishes at the optimal boundary of a smooth-fit solution, a mean absolute value is not an estimate
of $B$, and a scalar bias does not define a boundary displacement. We directly evaluate the splitting identity \eqref{eq:split} in this part.

For a stopping set $\cD_b$ on the grid, $P_b=G-J_b$ and $F_b$ solves the linear system
$(rI-\cL_h)F=\Gamma_h\ind_{\cC_b}$ on all nodes, where $\Gamma_h=(rI-\cL_h)G+c$ is the discrete gain.
\Cref{tab:bias} and \Cref{fig:bias} (right) report $B_b=F_b-P_b$ at $x_{\rm ref}$ for the computed
stopping set and for sets obtained by shifting both boundaries by $\delta$ (so that $\delta>0$
enlarges the continuation region). At the computed solution $B_b(x_{\rm ref})$ is $-0.046$, $-0.022$
and $-0.012$ for $N=121,241,481$: it halves with $h$, as expected from the $O(h)$ discrete
smooth-fit defect, and is consistent with a continuum value of zero. For shifted boundaries $B_b$ is
of order one, negative when the continuation region is too small and positive when it is too large,
exactly as in Example~\ref{ex:1d} and as predicted by Proposition~\ref{prop:defect}.

\begin{table}[htbp]
\begin{rv}{revA}{\tA{M4}}
\centering\small
\caption{Killed potential $P_b$, unrestricted potential $F_b$ and their difference $B_b$ at
$x_{\rm ref}$ for the computed stopping set ($\delta=0$) and for shifted boundaries.}\label{tab:bias}
\begin{tabular}{@{}rrrrr@{}}
\toprule
$N$ & shift $\delta$ & $P_b(x_{\rm ref})$ & $F_b(x_{\rm ref})$ & $B_b(x_{\rm ref})$\\
\midrule
$241$ & $-0.20$ & $-0.2266$ & $-1.2683$ & $-1.0417$\\
$241$ & $-0.10$ & $-0.2626$ & $-0.7527$ & $-0.4901$\\
$241$ & $-0.05$ & $-0.2723$ & $-0.5064$ & $-0.2341$\\
$241$ & $0.05$ & $-0.2703$ & $-0.0598$ & $0.2104$\\
$241$ & $0.10$ & $-0.2566$ & $0.1358$ & $0.3924$\\
$241$ & $0.20$ & $-0.1987$ & $0.4664$ & $0.6651$\\
\midrule
$121$ & $0$ & $-0.2752$ & $-0.3211$ & $-0.0459$\\
$241$ & $0$ & $-0.2752$ & $-0.2976$ & $-0.0224$\\
$481$ & $0$ & $-0.2752$ & $-0.2868$ & $-0.0116$\\
\bottomrule
\end{tabular}
\end{rv}
\end{table}

An independent Monte Carlo evaluation at the finite-difference boundary ($N=481$) distinguishes the
three quantities discussed after Proposition~\ref{prop:split}. With $Z$ the discretized post-stopping
functional accumulated from $\hat\tau$ to $T=120$ over visits to $\cC$ ($M=2\times10^4$,
$\Delta t=0.005$, $\varepsilon=\sqrt{\Delta t}$), we find
\[
\hat\E[Z]=0.0006\pm0.0081,\qquad\hat\E|Z|=0.455\pm0.005 ,
\]
with $95\%$ confidence half-widths. At the coarser step $\Delta t=0.01$ the values are
$0.0094\pm0.0079$ and $0.445\pm0.005$, so that a small positive discretization bias is visible in
the mean; extending the horizon to $T=200$ at $\Delta t=0.01$ changes these values by less than
$4\times10^{-4}$ (the discount factor $e^{-rT}$ is $2.5\times10^{-3}$ at $T=120$). The unrestricted process
does return to the continuation region after $\tau$, and the pathwise post-stopping functional is of
order one in absolute value, yet its mean, which estimates $B(x_{\rm ref})$, is zero within the statistical error at the finer
time step and of the order of the discretization bias at the coarser one. 
\end{rv}

\begin{rv}{revM}{\tA{Ed3, N4}\tB{23}}
\subsection{Comparative statics}

\Cref{tab:statics} varies $\alpha$, $r$, $c_0$ and $\rho$. These runs use $X_{\max}=4$ and $N=321$, i.e.\
the spacing $h=0.0125$ of the $N=241$ grid on a larger domain, so that the point of $\Da$ at which the
wedge is measured is well inside the domain also for $\alpha\ne1$. Increasing $r$ or $c_0$
increases $\Gac$, shrinks the continuation region and narrows the wedge. Changing $\alpha$ rotates
$\Da$ and makes the arms unequal. For the correlation sweep, $q$ decreases linearly from $0.576$ at
$\rho=-0.8$ to $0.064$ at $\rho=0.8$, and the wedge half-width and the premium on $\Da$ decrease
with it (\Cref{fig:rho}). The formal expressions of Remark~\ref{rem:width} reproduce the premium at
$(2,2)$ within $3\%$ over the sweep and the half-width within $3$--$16\%$ in the symmetric
cases; at $\rho=0.8$ the wedge spans only about two cells. For $\alpha\ne1$
the asymmetric formal half-widths \eqref{eq:width} are of similar accuracy (within $5$--$12\%$). We stress the interpretation:
$q$ is the density of the quadratic variation of $Y$, and it sets the size of the local-time reward
in \eqref{eq:gainform}. 

\begin{table}[htbp]
\begin{rv}{revM}{\tA{Ed3, N4}\tB{23}}
\centering\small
\caption{Comparative statics ($X_{\max}=4$, $N=321$, $h=0.0125$; areas over the window $[0,2.5]^2$). Half-widths $\delta_+$
(towards $\cD^{(1)}$) and $\delta_-$ (towards $\cD^{(2)}$) are measured in the variable
$y=x_1-\alpha x_2$ at the point of $\Da$ with $G=2$; the formal values are those of
\eqref{eq:width}.}\label{tab:statics}
\begin{tabular}{@{}lrrrrrrrrr@{}}
\toprule
 & & \multicolumn{3}{c}{area} & \multicolumn{2}{c}{computed} & \multicolumn{2}{c}{formal} & premium\\
\cmidrule(lr){3-5}\cmidrule(lr){6-7}\cmidrule(lr){8-9}
Case & $q$ & $\cC$ & $\cD^{(1)}$ & $\cD^{(2)}$ & $\delta_+$ & $\delta_-$ & $\delta_+$ & $\delta_-$ & on $\Da$\\
\midrule
$\alpha=0.7$ & $0.238$ & $2.74$ & $2.57$ & $1.00$ & $0.104$ & $0.094$ & $0.111$ & $0.100$ & $0.0243$\\
$\alpha=1.0$ (baseline) & $0.320$ & $2.66$ & $1.83$ & $1.83$ & $0.147$ & $0.147$ & $0.157$ & $0.157$ & $0.0398$\\
$\alpha=1.4$ & $0.474$ & $2.77$ & $1.02$ & $2.52$ & $0.233$ & $0.271$ & $0.254$ & $0.307$ & $0.0735$\\
$r=0.03$ & $0.320$ & $3.00$ & $1.65$ & $1.65$ & $0.159$ & $0.159$ & $0.170$ & $0.170$ & $0.0437$\\
$r=0.08$ & $0.320$ & $2.31$ & $2.00$ & $2.00$ & $0.130$ & $0.130$ & $0.140$ & $0.140$ & $0.0353$\\
$c_0=0.005$ & $0.320$ & $2.71$ & $1.80$ & $1.80$ & $0.148$ & $0.148$ & $0.158$ & $0.158$ & $0.0402$\\
$c_0=0.03$ & $0.320$ & $2.46$ & $1.92$ & $1.92$ & $0.141$ & $0.141$ & $0.151$ & $0.151$ & $0.0382$\\
$\rho=-0.8$ & $0.576$ & $3.12$ & $1.60$ & $1.60$ & $0.250$ & $0.250$ & $0.282$ & $0.282$ & $0.0719$\\
$\rho=-0.4$ & $0.448$ & $2.90$ & $1.71$ & $1.71$ & $0.200$ & $0.200$ & $0.220$ & $0.220$ & $0.0559$\\
$\rho=0.4$ & $0.192$ & $2.39$ & $1.96$ & $1.96$ & $0.090$ & $0.090$ & $0.094$ & $0.094$ & $0.0237$\\
$\rho=0.8$ & $0.064$ & $2.08$ & $2.12$ & $2.12$ & $0.026$ & $0.026$ & $0.031$ & $0.031$ & $0.0076$\\
\bottomrule
\end{tabular}
\end{rv}
\end{table}

\begin{figure}[htbp]
\begin{rv}{revA}{\tA{Ed3}}
\centering
\includegraphics[width=0.95\linewidth]{figs/rho.pdf}
\caption{Correlation sweep, $\rho\in\{-0.8,-0.6,\dots,0.8\}$ ($h=0.0125$). Left: wedge half-width at
$(2,2)$ against $q$, with the formal value $q/(4\Gac)$. Right: continuation premium $V-G$ at
$(2,2)$, with the formal value $q/(16\Gac)$.}\label{fig:rho}
\end{rv}
\end{figure}
\end{rv}

\begin{rv}{revM}{\tE{7}\tA{M6, N4}\tB{24}}
\subsection{Monte Carlo estimators of the killed potential}\label{sec:mcnum}

We evaluate the estimators \eqref{eq:gainest}--\eqref{eq:direst} at $x_{\rm ref}=(0.9,1.0)$ for the
finite-difference stopping set ($N=481$), whose killed potential is
$P(x_{\rm ref})=-(V-G)(x_{\rm ref})=-0.2752$ by \Cref{tab:refine}. Here $|Y_0|=0.1$.
\Cref{tab:mc} and \Cref{fig:mc} report the results for $M=2\times10^5$ paths, $T_{\max}=40$. The
mean killing time is between $3.6$ and $3.9$, depending on $\Delta t$, and the fraction of paths alive at $T_{\max}$ is below
$5\times10^{-5}$, so the truncation of the horizon is immaterial for the killed quantities, even
though $e^{-rT_{\max}}=0.135$ is not small.

\begin{table}[htbp]
\begin{rv}{revM}{\tA{M6, N4}\tB{24}}
\centering\small
\caption{Monte Carlo estimates of $P(x_{\rm ref})$ with $95\%$ confidence half-widths
($M=2\times10^5$; FD reference $-0.2752$). ``Paired'' is the mean of
$\hat P^{\rm gain}-\hat P^{\rm dir}$ on common paths for $\varepsilon=\sqrt{\Delta t}$. Confidence
intervals quantify sampling error only.}\label{tab:mc}
\begin{tabular}{@{}rcccccc@{}}
\toprule
 & & \multicolumn{4}{c}{gain-form estimator \eqref{eq:gainest}, bandwidth $\varepsilon=$} & paired\\
\cmidrule(lr){3-6}
$\Delta t$ & direct \eqref{eq:direst} & $\tfrac12\sqrt{\Delta t}$ & $\sqrt{\Delta t}$ & $2\sqrt{\Delta t}$ & $(\Delta t)^{0.4}$ & difference\\
\midrule
$0.02$ & $-0.2735\pm0.0011$ & $-0.2687$ & $-0.2711$ & $-0.2307$ & $-0.2531$ & $0.0024\pm0.0028$\\
$0.01$ & $-0.2742\pm0.0010$ & $-0.2723$ & $-0.2770$ & $-0.2549$ & $-0.2652$ & $-0.0028\pm0.0028$\\
$0.005$ & $-0.2754\pm0.0010$ & $-0.2732$ & $-0.2722$ & $-0.2678$ & $-0.2715$ & $0.0032\pm0.0028$\\
$0.0025$ & $-0.2754\pm0.0010$ & $-0.2744$ & $-0.2740$ & $-0.2742$ & $-0.2727$ & $0.0015\pm0.0028$\\
\bottomrule
\end{tabular}

\smallskip
{\footnotesize The half-widths of the gain-form estimates are between $\pm0.0017$ and $\pm0.0021$.}
\end{rv}
\end{table}

\begin{figure}[htbp]
\begin{rv}{revA}{\tA{M6, N4}}
\centering
\includegraphics[width=0.95\linewidth]{figs/mc.pdf}
\caption{Gain-form estimator at $x_{\rm ref}$. Left: estimate against bandwidth for four time steps
($M=2\times10^5$); the dashed line is the FD reference and the dotted line marks
$\varepsilon=|Y_0|$. Right: estimate against sample size ($\Delta t=0.005$,
$\varepsilon=\sqrt{\Delta t}$), with $95\%$ confidence intervals.}\label{fig:mc}
\end{rv}
\end{figure}

Three observations can be made. First, the direct estimator converges to the reference as $\Delta t$
decreases ($-0.2754\pm0.0010$ at $\Delta t\le0.005$), with a time-discretization bias of about
$0.002$ at $\Delta t=0.02$. Second, when $\varepsilon<|Y_0|=0.1$ the gain-form estimator differs from the direct one by at most
$0.005$, that is, by one to two confidence half-widths of the paired difference, and it is clearly
biased upwards when the band contains the starting level: at $\Delta t=0.01$ the estimates for $\varepsilon\in\{0.05,0.1,0.2\}$ are
$-0.2723$, $-0.2770$ and $-0.2549$, a spread of $0.022$, which is more than ten times the confidence
half-width. 
The estimator is sensitive to the bandwidth once the band reaches the
starting level, in line with Remark~\ref{rem:estimator}(d), and that in this experiment the bias is
governed by $\varepsilon$ itself, not by the ratio $\varepsilon/\sqrt{\Delta t}$. Third, the choice
$\varepsilon=(\Delta t)^{0.4}$, which is covered by Proposition~\ref{prop:consistency}, converges as
$\Delta t\to0$ but more slowly than $\varepsilon\propto\sqrt{\Delta t}$ in this range, because of
the smoothing bias. 
We therefore only make claims within the following ranges:
the bandwidth should be small compared with the distance of the starting state from $\Da$, the sensitivity to
$\varepsilon$ should be reported, and the direct estimator should be used as a control. The
half-widths in \Cref{tab:mc} scale as $M^{-1/2}$ (\Cref{fig:mc}, right).
\end{rv}

\begin{rv}{revM}{\tE{5}\tA{N1}}
\subsection{A run of the heuristic boundary iteration}\label{sec:algnum}

\Cref{alg:sf} was run on the baseline with nodes $s_k\in\{0,0.25,\dots,1.75\}$ for each arm,
$M=2\times10^4$ paths per trial state, $\Delta t=0.01$, $T_{\max}=25$, the gain-form estimator
\eqref{eq:gainest} with $\varepsilon=\tfrac12\sqrt{\Delta t}=0.05$, offset $h_n=0.05$, gains
$g_i=0.3/(1+i/10)$ with steps clipped to $\pm0.1$, margin $\varsigma=0.15$ and upper bounds
$\bar B_i=3.5$, for $I=60$ iterations, from two initial boundaries: a ``wide'' one,
$b_1(s)=b_2(s)=s+0.9$, and a ``narrow'' one, $b_1(s)=b_2(s)=\max(1,s+0.2)$. The two arms are updated
independently; symmetry is not imposed. \Cref{fig:alg} and \Cref{tab:alg} show the history. The
root-mean-square distance of the nodal values from the finite-difference boundary ($N=481$) falls
from $0.47$ and $0.37$ to about $0.01$ within ten iterations and then fluctuates; the average of
the last $15$ iterates is within $0.012$ of the finite-difference boundary at every node, and the
two runs agree with each other to within $0.01$. The remaining offset is systematic and of the size of the biases identified
in \Cref{sec:mcnum} (discretely monitored killing at $\Delta t=0.01$, finite offset $h_n$); it is not
reduced by averaging. The policies defined by the averaged iterates were then evaluated
independently by the finite-difference linear solve: their relative loss \eqref{eq:loss} has mean
$3\times10^{-5}$ and maximum $2.6\times10^{-3}$ over the window, the maximum occurring where the
parametrization continues the boundary parallel to $\Da$ beyond the last node.

\begin{figure}[htbp]
\begin{rv}{revM}{\tE{5}\tA{N1}}
\centering
\includegraphics[width=0.95\linewidth]{figs/alg.pdf}
\caption{\Cref{alg:sf} on the baseline. Left: root-mean-square distance of the nodal boundary values
from the FD boundary. Right: initial boundaries (dotted), averaged final iterates (markers; the two
runs nearly coincide) and the FD boundary (solid).}\label{fig:alg}
\end{rv}
\end{figure}

\begin{table}[htbp]
\begin{rv}{revM}{\tE{5}\tA{N1}}
\centering\small
\caption{History of \Cref{alg:sf} from two initial boundaries. Distances are between nodal values
and the FD boundary ($N=481$); ``defect'' is the root-mean-square of the estimated smooth-fit defects
$\hat d_i(s_k)$, whose standard error per node is about $0.06$; loss is the relative loss
\eqref{eq:loss} of the policy, evaluated by finite differences.}\label{tab:alg}
\begin{tabular}{@{}llcccc@{}}
\toprule
Run & iterate & r.m.s.\ distance & max.\ distance & r.m.s.\ defect & mean loss\\
\midrule
wide & initial & $0.472$ & $0.717$ & $3.54$ & $5.0\times10^{-2}$\\
 & $i=10$ & $0.015$ & --- & --- & ---\\
 & $i=60$ & $0.008$ & --- & $0.07$ & ---\\
 & average of last $15$ & $0.008$ & $0.011$ & --- & $3.0\times10^{-5}$\\
\midrule
narrow & initial & $0.368$ & $0.508$ & $1.10$ & $2.4\times10^{-2}$\\
 & $i=10$ & $0.012$ & --- & --- & ---\\
 & $i=60$ & $0.009$ & --- & $0.07$ & ---\\
 & average of last $15$ & $0.009$ & $0.012$ & --- & $2.9\times10^{-5}$\\
\bottomrule
\end{tabular}
\end{rv}
\end{table}

This is a single benchmark and two starting points; it shows that the smooth-fit iteration can work,
and that it selects a boundary through an optimality condition. It is not evidence of convergence in general. In particular the
estimated defect at the final iterates is at the level of its standard error, so the stopping
criterion is statistical, and the method inherits the discretization biases of the estimators.
\end{rv}
